\chapter{Definitions, Acronyms and Abbreviations}
\label{app:acronyms}

\section{Definitions}

The terms used in this report are defined in Table~\ref{tab:defs}.

\begin{table}[!htbp]
\centering
\caption[Definitions of terms used in the report]{Definitions of terms used in the report.}
\label{tab:defs}
\small
\begin{tabularx}{\textwidth}{@{}L{4.0cm}Y@{}}
\toprule
Term & Meaning in this report\\
\midrule
Assessable & A tooth record with an FDI code whose numbering confidence is at least 0.30\\
Constant (image-blind) & A fixed set of FDI codes answered for every image, without reading it; here the eight first and second molars\\
Fit set & The 78 DENTEX images, excluded from scoring because the models trained on them, on which the constant is fitted\\
Scored set & The 164 held-out DENTEX images on which the systems are compared\\
FDI-set F1 & Micro-averaged F1 over the set of tooth codes a system marks carious, per image\\
Lesion agreement & The share of expert lesions whose containing tooth is flagged by a system\\
Prior mass & The share of a cohort's lesions that lie on the constant's teeth\\
Matched-volume constant & An image-blind answer that names as many codes per image as a given system\\
Identity factors & Detection recall and numbering accuracy, the first two factors of the decision chain\\
Paired cluster bootstrap & Resampling radiographs (not teeth) with replacement, both systems on the same draw\\
Pre-registration & An acceptance rule committed to version control before the experiment it judges\\
Lever & A numbered, pre-registered proposal to change a pipeline component\\
Guard & A condition of a lever that must not get worse, such as Tufts detection accuracy\\
Abnormality flag & Stage 4a's output: the tooth crop overlaps an expert-marked Tufts abnormality\\
\bottomrule
\end{tabularx}
\end{table}

\section{Acronyms and Abbreviations}

The acronyms and abbreviations are expanded in Table~\ref{tab:acronyms}.

\begin{table}[!htbp]
\centering
\caption[Acronyms and abbreviations]{Acronyms and abbreviations.}
\label{tab:acronyms}
\small
\begin{tabularx}{\textwidth}{@{}L{2.2cm}YL{2.2cm}Y@{}}
\toprule
Acronym & Meaning & Acronym & Meaning\\
\midrule
AI & Artificial intelligence & JSON & JavaScript Object Notation\\
AP & Average precision & LLD & Low-level design\\
AP50 & Average precision at IoU 0.5 & mAP & Mean average precision\\
AUC & Area under the (ROC) curve & MD5 & Message-Digest 5 hash\\
CEJ & Cemento-enamel junction & NMS & Non-maximum suppression\\
CLAHE & Contrast-limited adaptive histogram equalisation & OPG & Orthopantomogram (panoramic radiograph)\\
CLI & Command-line interface & pHash, dHash & Perceptual and difference hashes\\
CNN & Convolutional neural network & PRS & Project requirements specification\\
CPU & Central processing unit & PRISMA & Preferred Reporting Items for Systematic Reviews and Meta-Analyses\\
DFD & Data-flow diagram & QUADAS-2 & Quality Assessment of Diagnostic Accuracy Studies, version 2\\
DL & DerSimonian--Laird (random effects) & ResNet & Residual network\\
ECE & Expected calibration error & ROC & Receiver operating characteristic\\
FDI & Fédération Dentaire Internationale (tooth notation) & RQ & Research question\\
FDP & Fixed dental prosthesis & SHA-256 & Secure Hash Algorithm, 256-bit\\
FPCL & Feature pyramid contrastive learning & UFBA-425 & The 425-radiograph dataset released with OralBBNet~\cite{budagam2024oralbbnet}\\
GPU & Graphics processing unit & VLM & Vision-language model\\
HKSJ & Hartung--Knapp--Sidik--Jonkman (interval) & YOLO & You Only Look Once (detector family)\\
HLD & High-level design & & \\
IoU & Intersection over union & & \\
\bottomrule
\end{tabularx}
\end{table}
